\documentclass[12pt,a4paper]{article}

\usepackage[utf8]{inputenc}
\usepackage[margin=1in]{geometry}
\usepackage{amsmath, amssymb, amsthm}
\usepackage{graphicx}
\usepackage{xcolor}
\usepackage{caption}
\captionsetup[figure]{font=small}
\usepackage{float}
\usepackage{subcaption}
\usepackage{booktabs}
\usepackage{tabularx}
\usepackage{siunitx}
\usepackage{hyperref}
\usepackage{bm}
%\usepackage{natbib}
\usepackage{multirow}
\usepackage{makecell}

% ---------- PATH MACROS FOR FIGURES ----------
% Seismic Amplitude plots
%\newcommand{\seismic}{../../../results/seismic_amplitude}
\newcommand{\figpath}{../figures}

\title{Seismic Amplitude Time series Analysis in the $H \times C$ Plane Using Ordinal Pattern Features}
\author{Rasika Dilhani}
\date{\today}

\begin{document}
	
	\maketitle
	
\section{Purpose}
The main purpose of this work is to investigate the temporal characteristics of seismic amplitude time series recorded at the WIZ.NZ seismic station using ordinal-pattern-based features. The study considers two types of seismic measurements, displacement and Real-time Seismic Amplitude Measurement (RSAM), derived from seismometer recordings sampled at 100 Hz and aggregated into 10-minute intervals from the year 2011 to 2022. Time series data were analysed at different temporal scales, including monthly datasets from January to December and quarterly datasets. The analysis employs ordinal pattern representations and information-theoretic measures to characterise the temporal behaviour of the seismic signals. Features derived from Shannon entropy, R\'enyi entropy, Tsallis entropy, statistical complexity, and Fisher information provide a comprehensive description of the underlying dynamics and enable the comparison of temporal patterns across different time periods. By examining these feature representations, the study aims to identify similarities and differences in seismic activity and assess the usefulness of ordinal-pattern-based features for analyzing and clustering seismic time series.
	
%Based on these features, we apply unsupervised time series clustering. The goal is to determine whether time series generated by different mechanisms can be grouped correctly using only their ordinal pattern features and $H \times C$ plane representations. The clustering results are evaluated using visual inspection of the $H \times C$ plane and silhouette coefficients from the $k$-means algorithm.
	
	
\section{Methodology}
\label{sec:methodology}

This study analyses seismic amplitude time series recorded at the WIZ.NZ seismic station between 2011 and 2022. Two types of measurements are considered: displacement and Real-time Seismic Amplitude Measurement (RSAM). The original seismometer recordings were sampled at 100~Hz and aggregated into 10-minute intervals. The analysis is conducted using both monthly and quarterly time series to investigate temporal variations in seismic activity at different time scales.

\subsection{Feature Construction}

Ordinal-pattern-based features are extracted from each time series according to the following procedure:

\begin{itemize}
	
	\item \textbf{Embedding:}
	A fixed embedding dimension ($D=5$) and time delay ($\tau=1$) are used to construct ordinal patterns from the seismic time series.
	
	\item \textbf{Temporal Aggregation:}
	The analysis considers monthly and quarterly subsets of the displacement and RSAM time series for the period 2011--2022.
	
	\item \textbf{Feature Extraction:}
	For each time series, the ordinal pattern distribution is obtained and the following quantities are computed:
	
	\begin{itemize}
		
		\item \emph{Entropy Measures:}
		Shannon entropy, R\'enyi entropy, Tsallis entropy, and Fisher information measure.
		
		\item \emph{Complexity Measures:}
		The corresponding Shannon, R\'enyi, Tsallis, and Fisher complexity measures.
		
		\item \emph{Disequilibrium:}
		The Jensen--Shannon disequilibrium associated with the ordinal pattern distribution.
		
		\item \emph{Variance Measures:}
		Variance estimates for selected entropy and complexity measures to assess feature stability.
		
		\item \emph{Confidence Intervals:}
		Confidence intervals for entropy and complexity measures are computed to quantify estimation uncertainty.
		
	\end{itemize}
	
\end{itemize}

The analysis is based on the following thirteen ordinal-pattern features:

\begin{itemize}
	\item $H_{\text{Shannon}}$,
	$H_{\text{Rényi}}$,
	$H_{\text{Tsallis}}$,
	$H_{\text{Fisher}}$,
	
	\item $C_{\text{Shannon}}$,
	$C_{\text{Rényi}}$,
	$C_{\text{Tsallis}}$,
	$C_{\text{Fisher}}$,
	
	\item Disequilibrium,
	
	\item $\mathrm{Var}(H_{\text{Shannon}})$,
	$\mathrm{Var}(H_{\text{Rényi}})$,
	$\mathrm{Var}(H_{\text{Tsallis}})$,
	
	\item $\mathrm{Var}(C_{\text{Shannon}})$.
\end{itemize}

\subsection{Information Theoretic Quantifiers}

For each ordinal pattern probability vector $\mathbf P=\{\mathbf p_\ell\}_{\ell=1}^{k}$, we evaluate a collection of 
global and local information measures.  
Given an estimator $\hat{\mathbf P}=\{\hat{\mathbf p}_\ell\}$ obtained from the empirical ordinal pattern frequencies, our goal is to study the behaviour of the entropy functionals $H(\hat{ \mathbf P})$ and their associated complexity measures.

\paragraph{Shannon entropy.}
The Shannon entropy is defined as
\[
H_{\mathrm{S}}(\hat{\mathbf P})
= -\sum_{\ell=1}^{k} \hat{\mathbf p}_\ell \log \hat{\mathbf p}_\ell,
\]
and we use its normalized version $H_{\mathrm{S}}/H_{\max}$, where $H_{\max}$ corresponds to the 
uniform distribution.

\paragraph{Tsallis entropy.}
For $q \in \mathbb{R}\setminus\{1\}$, the Tsallis entropy is given by
\[
H^{\mathrm{T}}_{q}(\hat{\mathbf P})
= \frac{1}{q-1}\left(1 - \sum_{\ell=1}^{k} \hat{\mathbf p}_\ell^{\,q}\right),
\]
which reduces to Shannon entropy in the limit $q \to 1$.  
A normalized version is obtained by dividing by its value under the uniform distribution.

\paragraph{Rényi entropy.}
For $q>0$, $q\neq 1$, the Rényi entropy is defined as
\[
H^{\mathrm{R}}_{q}(\hat{P})
= \frac{1}{1-q}\log\!\left(\sum_{\ell=1}^{k} \hat{\mathbf p}_\ell^{\,q}\right),
\]
and is likewise normalized using its maximum value at the uniform distribution.

\paragraph{Fisher Information Measure.}
To quantify local sensitivity in the ordinal pattern distribution and this is discussed by Felipe et.al.~\cite{Olivares2012b}, we adopt the discrete Fisher Information Measure proposed by Ferri et.al.~\cite{Ferri2009}:
\[
F(\hat{\mathbf P})
= \frac{1}{2}\sum_{\ell=1}^{k-1}
\frac{(\hat{\mathbf p}_{\ell+1}-\hat{\mathbf p}_{\ell})^{2}}
{\hat{\mathbf p}_{\ell+1}+\hat{\mathbf p}_{\ell}}.
\]

This expression emphasizes differences between adjacent probabilities and therefore captures 
fine scale structure that entropy based measures may overlook.

\subsection{Data Representation}\label{data-representation}
The seismic dataset consists of temporal observations recorded at the WIZ seismic station in New Zealand. Two measures of volcanic seismic activity were analysed: the Real-time Seismic Amplitude Measurement (RSAM) and the average ground displacement. The data cover the period from 2011 to 2022 and were obtained from continuous seismic monitoring records. Each observation is associated with a timestamp, allowing the construction of long-term time series that capture the temporal evolution of seismic activity. These time series were used to investigate temporal patterns, variability, and characteristic behaviours of seismic signals through ordinal-pattern-based information-theoretic descriptors.

The following figures gives the time series plots generated from annual displacement and RSAM, mean monthly displacement and monthly RSAM.

\begin{figure}[H]
	\centering
	\includegraphics[width=\textwidth]{\figpath/annual_displacement}
	\caption{Annual Displacement.}
	\label{fig:annualdisplacement}
\end{figure}

\begin{figure}[H]
	\centering
	\includegraphics[width=\textwidth]{\figpath/annual_rsam}
	\caption{Annual RSAM.}
	\label{fig:annualrsam}
\end{figure}

\begin{figure}[H]
	\centering
	\includegraphics[width=\textwidth]{\figpath/all_in_one_RSAM_by_year}
	\caption{All in one Annual RSAM.}
	\label{fig:all_in_one_RSAM}
\end{figure}

\begin{figure}[H]
	\centering
	\includegraphics[width=\textwidth]{\figpath/monthly_mean_displacement}
	\caption{Mean monthly Displacement.}
	\label{fig:mean_displacement}
\end{figure}

\begin{figure}[H]
	\centering
	\includegraphics[width=\textwidth]{\figpath/monthly_mean_RSAM}
	\caption{Mean monthly RSAM.}
	\label{fig:mean_RSAM}
\end{figure}

\subsection{Proposed Methodology}\label{proposed-methodology}

The proposed methodology begins with the collection and preprocessing of seismic time series data, including average ground displacement and Real-time Seismic Amplitude Measurement (RSAM) observations recorded at the WIZ seismic station. The data are organised into annual time series spanning the period 2011-2022 and transformed into a common temporal framework to facilitate comparisons across years. Exploratory analyses are then conducted through visualisation of the temporal behaviour of the displacement and RSAM series, including annual, monthly, and multi-year representations. These analyses provide an initial assessment of the variability, trends, and temporal characteristics of the seismic signals and establish the foundation for subsequent ordinal-pattern-based feature extraction and information-theoretic analysis.



\bibliographystyle{acm}
%\bibliography{../../BearingFaultDiagnosis}
\bibliography{../../Ordinal_Patterns_R/Text/BearingFaultDiagnosis}

\end{document}